\documentclass[conference]{IEEEtran}

\usepackage{amsmath}
\usepackage{cite}
\usepackage{url}

\begin{document}

\title{Milestone 3: Critical Evaluation and Transfer\\
Loop Unrolling and Controller Delay in High-Level Synthesis}

\author{\IEEEauthorblockN{Christian Percival}
\IEEEauthorblockA{Hochschule Hamm-Lippstadt\\
Electronic Engineering}}

\maketitle

\section{Milestone 3: Critical Evaluation and Transfer}

\subsection{Strengths of the Approach}

\begin{itemize}
    \item The paper identifies an important HLS problem that can easily be overlooked.
    \item It shows that loop unrolling affects not only latency, but also controller delay.
    \item It connects a high-level compiler optimization to low-level hardware timing effects.
    \item The approach avoids exhaustive synthesis of all unroll factors.
    \item The estimation method can reduce exploration time significantly.
    \item The priority-based selection method gives a practical way to choose promising unroll factors.
    \item The paper supports the idea that HLS optimization must consider hardware constraints.
\end{itemize}

\subsection{Limitations of the Approach}

\begin{itemize}
    \item The delay model depends on the synthesis flow and target technology.
    \item The original experiments use a specific behavioral synthesis framework and ASIC library.
    \item The exact delay constants may not apply directly to FPGA-based HLS tools.
    \item Modern HLS tools may optimize the design differently.
    \item Controller delay may be difficult to isolate in modern synthesis reports.
    \item The method focuses mainly on controller delay, while real timing also includes datapath, routing, memory, and placement effects.
    \item The approach may be less accurate for designs with complex memory behavior or tool-specific optimizations.
\end{itemize}

\subsection{Hidden or Restrictive Assumptions}

\begin{itemize}
    \item The loop structure must be suitable for unrolling.
    \item Loop bounds should be known or predictable.
    \item The effect of unrolling should be estimable from high-level information.
    \item The controller delay model must be calibrated for the target synthesis environment.
    \item The approach assumes that delay estimation is accurate enough to guide design-space exploration.
    \item It also assumes that reducing exploration time is more important than finding an absolutely optimal result through exhaustive synthesis.
\end{itemize}

\subsection{Risks}

\begin{itemize}
    \item An inaccurate delay estimate could select an unroll factor that fails timing.
    \item A conservative estimate could reject useful unroll factors.
    \item Modern HLS tools may hide controller details, making validation difficult.
    \item A student implementation may only measure total timing and resource usage, not pure controller delay.
    \item If the benchmark is too simple, the controller-delay effect may not be very visible.
    \item If the benchmark is too complex, the result may become difficult to explain.
\end{itemize}

\subsection{Suitable Application Examples}

\begin{itemize}
    \item The approach is suitable for loop-based hardware accelerators.
    \item Suitable examples include:
    \begin{itemize}
        \item dot-product accelerators,
        \item FIR filters,
        \item simple image-processing kernels,
        \item matrix or vector operations,
        \item embedded signal-processing loops.
    \end{itemize}

    \item These applications often contain repeated loop operations and can benefit from partial unrolling.
    \item The approach is especially useful when there is a strict timing budget.
\end{itemize}

\subsection{Unsuitable or Problematic Application Examples}

\begin{itemize}
    \item The approach may be less suitable for code with highly irregular control flow.
    \item It may be less useful when loop bounds are unknown or data-dependent.
    \item It may be difficult to apply when memory access patterns dominate performance.
    \item It may be less effective if the HLS tool already performs aggressive automatic optimization.
    \item It may also be less useful for very small loops where exhaustive exploration is already cheap.
\end{itemize}

\subsection{Transfer to a Student Implementation}

\begin{itemize}
    \item The original ASIC-based synthesis setup is likely not available for a low-budget student project.
    \item A practical student implementation can instead use:
    \begin{itemize}
        \item a Python-based estimation model,
        \item AMD Vitis HLS,
        \item Intel HLS,
        \item Bambu HLS,
        \item or another available university HLS tool.
    \end{itemize}

    \item The implementation should not try to reproduce the original paper exactly.
    \item The goal should be to demonstrate the same engineering principle.
    \item The student experiment can compare different unroll factors for a small benchmark loop.
    \item Example unroll factors:
    \begin{itemize}
        \item 1,
        \item 2,
        \item 4,
        \item 8,
        \item 16,
        \item 32.
    \end{itemize}

    \item For each unroll factor, the following should be recorded:
    \begin{itemize}
        \item latency in cycles,
        \item estimated clock period,
        \item timing slack,
        \item LUT or ALM usage,
        \item flip-flop or register usage,
        \item DSP usage,
        \item BRAM or memory usage,
        \item synthesis time if available.
    \end{itemize}
\end{itemize}

\subsection{Possible Python Simulation}

\begin{itemize}
    \item A Python model can be used if HLS tools are not available.
    \item The Python model can estimate:
    \begin{itemize}
        \item operation count,
        \item number of loop iterations after unrolling,
        \item estimated latency,
        \item estimated number of states,
        \item estimated state bits,
        \item estimated controller delay.
    \end{itemize}

    \item The simulation would not prove real hardware timing.
    \item However, it can help demonstrate the trend that unrolling reduces loop iterations while increasing controller complexity.
\end{itemize}

\subsection{Possible HLS Experiment}

\begin{itemize}
    \item A small dot-product function can be written in C or C++.
    \item The loop can be synthesized multiple times with different unroll factors.
    \item The results can be compared using the HLS synthesis reports.
    \item The most important output would be a table and graph showing:
    \begin{itemize}
        \item unroll factor versus latency,
        \item unroll factor versus resource usage,
        \item unroll factor versus clock estimate or timing slack.
    \end{itemize}

    \item This would provide a practical demonstration of the paper's main idea.
\end{itemize}

\subsection{Possible Extensions}

\begin{itemize}
    \item Compare loop unrolling with loop pipelining.
    \item Study how memory partitioning affects the benefit of loop unrolling.
    \item Test more than one benchmark loop.
    \item Compare full unrolling with partial unrolling.
    \item Extend the model to include datapath delay and routing delay.
    \item Investigate how modern FPGA HLS tools report controller-related timing.
\end{itemize}

\subsection{Critical Conclusion}

\begin{itemize}
    \item The paper provides a useful reminder that HLS optimizations must be evaluated from a hardware perspective.
    \item Loop unrolling can reduce latency, but it can also increase controller complexity and timing pressure.
    \item The best unroll factor is usually a tradeoff, not simply the largest possible factor.
    \item Estimation-based exploration is useful because exhaustive synthesis can be too slow.
    \item For a student project, the most realistic implementation is a simplified estimation model supported by a small HLS experiment if tools are available.
    \item The final seminar paper should clearly distinguish between:
    \begin{itemize}
        \item what the original paper proves,
        \item what can be reproduced in a student environment,
        \item and what remains a limitation or open question.
    \end{itemize}
\end{itemize}

\begin{thebibliography}{00}

\bibitem{kurra2007}
S. Kurra, N. K. Singh, and P. R. Panda, ``The Impact of Loop Unrolling on Controller Delay in High Level Synthesis,'' DATE, 2007.

\end{thebibliography}

\end{document}
